\section{Method}
\label{sec:method}

% Method overview (TikZ). Palette: peach #F7D8B8, light aqua #BFE6EA, salmon #F6A69A, mist #D9E2EB, warm white #F2F2F2,
% deep teal #1F8FAE, coral #E8685B, dark coral #A83E35.
\begin{figure*}[t]
\centering
\resizebox{\textwidth}{!}{%
\begin{tikzpicture}[
  font=\footnotesize,
  >={Stealth[length=2.2mm]},
  mod/.style={draw=#1!70!black, fill=#1, rounded corners=3pt, minimum height=1.35cm, text width=2.7cm, align=center, inner sep=3pt},
  io/.style={draw=gray!60, fill=cWarm, rounded corners=2pt, align=left, inner sep=3pt, text width=2.3cm},
  lab/.style={font=\scriptsize\itshape, text=gray!80!black},
  flow/.style={->, thick, draw=gray!70!black},
  save/.style={->, thick, dashed, draw=cCoral},
]
% input
\node[io] (doc) at (0,0) {\textbf{Long document $x$}\\[1pt]
  \tiny \ldots{} the special magic uuid for dusty-harbor is: \colorbox{cSalmon}{\texttt{186c24bd-2192-4b86-\ldots}} \ldots{} founders in 1995 \ldots{}};
% modules
\node[mod=cPeach, right=0.55cm of doc] (det) {\textbf{\textsc{Detect}}\\[1pt]\scriptsize text-side hazard set $\mathcal{H}(x)$\\ regex + entropy; years excluded};
\node[mod=cAqua, right=0.55cm of det] (alloc) {\textbf{\textsc{Allocate}}\\[1pt]\scriptsize prose at $s_{\mathrm{lo}}$, hazards at $s_{\mathrm{hi}}$\\ $\max s_{\mathrm{hi}}$ s.t.\ $v(I)\le B$};
\node[mod=cSalmon, right=0.55cm of alloc] (anc) {\textbf{\textsc{Anchor}}\\[1pt]\scriptsize verbatim anchors of $\mathcal{H}(x)$, paid by $B - v(I)$};
\node[mod=cMist, right=0.55cm of anc] (route) {\textbf{\textsc{Route}}\\[1pt]\scriptsize $r=\frac{\text{anchor tokens}}{v(I)}$\\ $r\le\tau$: image; else text};
% outputs
\node[io, text width=1.9cm, above right=0.05cm and 0.5cm of route.east, anchor=south west] (img) {\textbf{VLM input}\\ \scriptsize image pages + anchors};
\node[io, text width=1.9cm, below right=0.05cm and 0.5cm of route.east, anchor=north west] (txt) {\textbf{VLM input}\\ \scriptsize plain text (hazard-dense)};
% flows
\draw[flow] (doc) -- (det);
\draw[flow] (det) -- node[lab, above] {$\mathcal{H}$} (alloc);
\draw[flow] (alloc) -- node[lab, above] {$I$} (anc);
\draw[flow] (anc) -- (route);
\draw[flow] (route.east) -- ++(0.2,0) |- (img.west);
\draw[flow] (route.east) -- ++(0.2,0) |- (txt.west);
\draw[save] (alloc.south) to[bend right=25] node[lab, below, text=cCoral] {saved budget $B-v(I)$} (anc.south);
\end{tikzpicture}}
\caption{\textbf{Hazard-aware rendering.} \textsc{Detect} marks verbatim-risky strings in the text before rendering; \textsc{Allocate} enlarges only those strings and shrinks prose so that the image never exceeds the vision-token budget $B$ of uniform rendering (a 6\,px uniform image in our experiments); \textsc{Anchor} spends the budget the allocation saved on anchors, verbatim text copies of the flagged strings; \textsc{Route} sends documents whose anchors would be too expensive to plain text. Every step is training-free and needs no access to the model.}
\label{fig:method}
\end{figure*}


\paragraph{Overview.}
Given a long document, an optical compressor renders it into image pages that a VLM reads in place of the text.
Our goal is to make the strings that must be copied exactly survive this rendering, without spending more vision tokens and without access to the model.
\Cref{fig:method} shows the pipeline.
\textsc{Detect} reads the text and marks verbatim-hazardous strings.
\textsc{Allocate}, our contribution, renders those strings large and the remaining prose small, so that the image costs no more than uniform rendering.
The other two steps follow prior work on failure-aware optical compression~\cite{anon2026renderretain}.
\textsc{Anchor} appends verbatim text copies of the flagged strings, paid for here by the vision tokens that allocation saves.
\textsc{Route} falls back to plain text when the anchors would be too expensive.

\subsection{Problem formulation}
\label{sec:formulation}

A document $x=(w_1,\dots,w_n)$ is a sequence of whitespace-separated words.
A renderer $R$ maps $x$ and a size assignment $\sigma:\{1,\dots,n\}\to\mathbb{N}$ (glyph height in pixels for every word) to an image $I=R(x;\sigma)$ of fixed width, which is split into pages for the VLM.
To compare renderings, we need their cost in the unit providers bill.
Following the patch-based accounting of dynamic-resolution VLMs~\cite{wang2024qwen2vl}, we define the vision-token cost of an image with pages of width $w_p$ and height $h_p$ as
\begin{equation}
  v(I) \;=\; \sum_{p} \left\lceil \tfrac{w_p}{28} \right\rceil \left\lceil \tfrac{h_p}{28} \right\rceil ,
  \label{eq:vtok}
\end{equation}
where each vision token covers a $28\times28$-pixel area (14-pixel patches merged $2\times2$ in Qwen2-VL).
\Cref{eq:vtok} is used only to enforce budgets; all costs we report are the prompt tokens billed by each provider.
The \emph{budget} is the cost of uniform rendering at a base size $s_0$, $B = v\big(R(x;\sigma\equiv s_0)\big)$.
Every method we compare uses the same number of vision tokens as the uniform baseline, up to page rounding.
Let $\mathcal{H}(x)\subseteq\{1,\dots,n\}$ be the indices of hazardous words, those whose value is needed verbatim downstream.
We seek a size assignment that makes every $w_i$, $i\in\mathcal{H}(x)$, as legible as possible subject to $v(I)\le B$, without deleting or reordering any character.

\subsection{Where pixels matter}
\label{sec:principle}

Which words should receive the pixels?
Consider how a reader copies a string $c_1\dots c_L$ from the image.
As a schematic model, the reader combines two terms for each character: visual evidence from the glyph, which grows with the glyph size $s$, and the expectation of its language model given the context:
\begin{equation}
  p(c_i \mid I, \mathrm{ctx}) \;\propto\; \underbrace{p(I_i \mid c_i;\, s_i)}_{\text{pixels}} \; \underbrace{p_\theta(c_i \mid \mathrm{ctx})}_{\text{language prior}} .
  \label{eq:reading}
\end{equation}
The pixels must supply whatever the prior leaves open, roughly the \emph{residual entropy} $H(c_i \mid \mathrm{ctx})$ of each character.
For English prose this is small, because most letters are predictable from the words around them.
For the hyphens of a UUID it is close to zero.
For its hexadecimal characters it is 4 bits, all of which must come from the glyph.
Exact copying fails if any single character fails, so a string is only as legible as its least predictable character.
Uniform rendering gives every character the same visual evidence, which is more than prose needs and, at dense budgets, less than an identifier needs.
When the evidence for a high-entropy character is too weak, the prior dominates \cref{eq:reading}. The prior knows the format but not the content, so we expect the reader to emit a well-formed string with invented characters.
Under a fixed budget, pixels should therefore move from characters with low residual entropy to characters with high residual entropy.
The trade is favorable because the two kinds of characters sit on different parts of their legibility curves.
Predictable text is on the flat part: in our controlled study, found English phrases keep 92--96\% character accuracy from 10 down to 4\,px, so taking pixels from prose costs little.
High-entropy strings at dense budgets are on the steep part: giving a UUID 12 instead of 5--6\,px raises its hexadecimal accuracy from 51\% to 91\% (\cref{sec:why}).
A small loss where the prior compensates buys a large gain where it cannot.
Under the patch convention of \cref{eq:vtok}, a 6\,px Menlo character is 4 pixels wide. One 14-pixel patch then mixes about 3.5 characters, and one vision token covers 7 characters on each of 3 lines; at 12\,px a patch holds 2 characters.
The two steps below implement this principle with simple means: a text-side proxy for residual entropy (\textsc{Detect}) and a two-level size assignment (\textsc{Allocate}).

\subsection{\textsc{Detect}: text-side hazard detection}
\label{sec:detect}

The detector approximates $\mathcal{H}(x)$ without labels, following the regex-plus-windowed-entropy design of~\cite{anon2026renderretain}.
Its premise is that a string is hazardous when it is \emph{high-entropy}: a misread character yields another valid-looking string, and the language model cannot repair it from context.
After stripping surrounding punctuation, we flag a word if it matches one of five patterns: a hexadecimal run of at least 8 characters, a run of at least 5 digits, a path containing a digit or a file extension, an IP address, or a digest prefix.
We also flag words that mix letters and digits and either have at least 20 identifier characters or contain a 16-character window with Shannon entropy of at least 3.5 bits.
Two choices encode what \emph{not} to protect.
Four-digit numbers are excluded because in natural text they are mostly years, which are low-entropy and cheap to read.
The length and entropy rules require mixed letter--digit tokens so that long English words are not flagged.
These exclusions are our changes.
A preliminary detector without them flagged 85--112 words per LongBench document on average, most of them years, and shrank prose in documents that did not need it (\cref{app:negative}).

When the detector flags nothing, the rendering is identical to the uniform baseline, so a document whose hazards all go undetected loses only the potential gain (\cref{sec:niq}).
A missed hazard in a document where \emph{other} strings are flagged, however, shrinks together with the prose.

\subsection{\textsc{Allocate}: budget-constrained glyph allocation}
\label{sec:allocate}

Given $\mathcal{H}(x)$, we restrict size assignments to two levels: prose words at $s_{\mathrm{lo}}$ and flagged words at $s_{\mathrm{hi}}$.
Prose is wrapped at character boundaries, as in the uniform baseline; each flagged word is kept whole on one line.
We choose the sizes by
\begin{equation}
\begin{aligned}
  (s_{\mathrm{hi}}^{\ast}, s_{\mathrm{lo}}^{\ast}) = &\operatorname*{lexmax}_{\substack{s_{\mathrm{hi}}\in[s_0,\,12]\\ s_{\mathrm{lo}}\in[5,\,\min(s_{\mathrm{hi}},s_0)]}} (s_{\mathrm{hi}}, s_{\mathrm{lo}})\\
  &\text{s.t.}\;\; v\big(R(x;\sigma_{s_{\mathrm{lo}},s_{\mathrm{hi}}})\big)\le B ,
\end{aligned}
  \label{eq:alloc}
\end{equation}
where $\sigma_{s_{\mathrm{lo}},s_{\mathrm{hi}}}$ assigns $s_{\mathrm{hi}}$ to words in $\mathcal{H}(x)$ and $s_{\mathrm{lo}}$ to all others, and $\operatorname{lexmax}$ first maximizes $s_{\mathrm{hi}}$ and then $s_{\mathrm{lo}}$.
The order follows \cref{sec:principle}: the marginal pixel is worth more on a hash than on prose.
The lower bound of 5\,px keeps prose within the range VLMs still read, and $s_{\mathrm{lo}}\le s_0$ ensures prose is never enlarged.
If no pair with $s_{\mathrm{hi}}>s_{\mathrm{lo}}$ fits, prose-only lines lose one pixel of line spacing and the search is repeated; at $s_0=5$\,px this is what makes room (\cref{app:alloc}).
If $\mathcal{H}(x)=\emptyset$, the result is identical to uniform rendering.
Allocation takes about 1\,s for 32K and 4\,s for 128K characters on one CPU core.

\subsection{\textsc{Anchor} and \textsc{Route}}
\label{sec:anchor}
\label{sec:route}

When prose can shrink ($s_0=6$\,px), the allocated image usually ends \emph{below} the budget, because the search stops at the largest feasible $s_{\mathrm{hi}}$ rather than filling every remaining row.
In our long-context runs, it used 13--27\% fewer billed tokens than uniform rendering.
Following~\cite{anon2026renderretain}, we spend these savings on \emph{anchors}: after the question, the flagged strings are appended verbatim as text, under a line stating that they are copied exactly from the image.
Anchors carry no key, but a document with only one or a few flagged strings yields a short list that may consist of the answer alone.
A reader can then answer without reading the image (\cref{sec:anchors}).
Anchor cost grows linearly with the number of flagged strings, so before any model call we compute the ratio
\begin{equation}
  r(x) \;=\; \frac{\sum_{i\in\mathcal{H}(x)} c(w_i)}{v(I)}, \qquad c(w) = \tfrac{|w|}{2} + 3,
  \label{eq:route}
\end{equation}
where $c(w)$ estimates the text tokens of one anchor and $I$ is the allocated image, and send the document as plain text when $r(x)>\tau$.
This replaces the density test of~\cite{anon2026renderretain} with one that accounts for the allocated image.
We use $\tau=0.3$, the median of per-reader two-fold cross-validation on a synthetic density sweep (\cref{app:density}).
